%% =====================================================================
%%  MMCS Group Project --- Tier 1 Technical Report --- TEMPLATE
%%  Compile:  tectonic -X compile MMCS_TechnicalReport_template.tex
%%
%%  This file is the authoring skeleton. Grey italic text in square
%%  brackets is guidance: delete it as you write. Numbers shown in the
%%  example tables are illustrative, not the group's results -- replace
%%  them, keep the column structure.
%%
%%  Page budget (body): 1 title+summary | 1 boundary/data | 2 model |
%%  3 results | 2 sensitivity | 1 recommendation | 1 validation/appendix.
%%  Hard ceiling: 15 pages including appendix. Longer is not better.
%%
%%  Two rules this file encodes, learned from past reports:
%%   1. every table has exactly one flexible column, so nothing is
%%      stretched to fill the page;
%%   2. numeric columns use siunitx S specs and contain numbers only --
%%      a dash or N/A goes in braces {N/A}.
%% =====================================================================
\documentclass[11pt,a4paper]{article}
\usepackage{mmcsreport}
\usepackage{multicol}

\renewcommand{\reportfootertext}{Confidential --- prepared for course assessment}
\hypersetup{pdftitle={MMCS Group Project -- Tier 1 Technical Report (template)},
            pdfauthor={Group 15}}

% units row under a table header: \unitcell is a font switch, not a command
\newcommand{\unitcell}{\footnotesize\sffamily\color{muted}}

% placeholder for a figure slot of a given height
\newcommand{\figslot}[2]{%
  \begin{tcolorbox}[enhanced,boxrule=0.4pt,colframe=rule,colback=white,arc=0pt,
    height=#1,valign=center,halign=center,fontupper=\footnotesize\color{muted}]
    #2\end{tcolorbox}}

\begin{document}
\thispagestyle{firststyle}

% ---------------------------------------------------------------- title ----
\eyebrow{MMCS Group Project \textperiodcentered{} 2026 \textperiodcentered{} Tier 1}
\reporttitle{Recovering the value of surplus\\renewable generation}
\reportsubtitle{Battery or electrolyser, at what scale, and under which contract}

\begin{reportmeta}
  \reportmetaitem{Prepared by}{Group 15}
  \reportmetaitem{Version}{1.0 --- 5 November 2026}
  \reportmetaitem{Data window}{1 Jan 2020 -- 17 Sep 2026}
  \reportmetaitem{Decision owner}{Portfolio board}
\end{reportmeta}

% ------------------------------------------------------------ key figures --
% Four numbers an examiner may read and stop. Choose the four that carry the
% recommendation: value created, the configuration, the binding constraint,
% and the payback. Do not put model statistics here.
\begin{keyfigures}[Headline results]
  \kpi{+£8.57m}{annual value of the recommended 20 MW module}
  \kpi{20 MW}{electrolyser; no battery is economic}
  \kpi{2.6 yr}{payback before stack replacement}
  \kpi{0.07\%}{of current sales diverted to make hydrogen}
\end{keyfigures}

% --------------------------------------------------------------- decision --
% One box, at most five lines. Cost ceiling -> revenue -> payback -> the
% condition under which the answer changes. This is the paragraph a busy
% reader quotes.
\begin{decisionbox}[Decision]
\ph{Build / do not build — one sentence. Then: what it costs, what it earns,
what it pays back, and the single condition that would reverse the answer.}
\end{decisionbox}

% ============================================================ section 1 ====
\section{Executive summary}

\ph{First paragraph contains numbers, not framing. Capacity, annual value,
payback, and the assumption that carries the answer. Two sentences of context
at most.}

\ph{Second paragraph: the mechanism. Why this configuration and not the
alternative, in one or two sentences, naming the constraint that actually
binds (a contract cap, a cost threshold, an operating rule) rather than the
one a reader would guess.}

\ph{Third paragraph: what the reader must accept. Data vintage, time
resolution, foresight, and the one number that would move most if that
assumption failed.}

\begin{caveatbox}[The assumption that carries the answer]
\ph{Name it and quantify it. ``The result holds while the offtake contract is
at least X; at Y the investment is zero.'' A caveat box that cannot be
falsified by a number is decoration.}
\end{caveatbox}

% ============================================================ section 2 ====
\section{Problem, boundary and data}

\subsection{The question}

\ph{State the decision in one sentence, then the arithmetic that frames it:
how much energy the system produces, how much demand it serves, how much is
stranded, and on how many days. These four numbers justify the whole report.}

\needspace{7\baselineskip}
\subsection{System boundary}

\ph{What is inside the model and what is not. Be explicit that constructed
supply is not observed curtailment, and that the demand series is the
absorption cap. One sentence per excluded mechanism (network, intraday,
ancillary, storage of hydrogen), because each exclusion is later a
limitation.}

\figslot{4.6cm}{Figure slot 1 --- system boundary diagram, 16.4\,cm $\times$ 4.6\,cm\\
\footnotesize generation $\rightarrow$ \{sale, battery, electrolyser, unused\} $\rightarrow$ demand}

\subsection{Data and provenance}

% Provenance table: every series, its period, its repair status, its use.
% This table is what separates a report that can be audited from one that
% cannot. Add a release tag or a hash in the source note.
\begin{table}[H]
\caption{Data series, provenance and use.}
\small
\begin{tabularx}{\linewidth}{@{}L{3.9cm}L{2.2cm}YL{3.0cm}@{}}
\toprule
\sffamily\bfseries Series & \sffamily\bfseries Period & \sffamily\bfseries Treatment & \sffamily\bfseries Where used \\
\midrule
\ph{Generation by fuel} & \ph{2020--2026} & \ph{repaired release v2; 523 of 2,452 days carry a non-strict flag} & \ph{all cases} \\
\ph{Demand} & \ph{2020--2026} & \ph{truncated last day of each query restored} & \ph{absorption cap} \\
\ph{Market index price} & \ph{2020--2026} & \ph{half-scale placeholder providers removed} & \ph{revenue} \\
\bottomrule
\end{tabularx}
\srcnote{Source: \ph{repository path and release tag}. Repair log: \ph{file}.}
\end{table}

\subsection{Conventions and assumptions}

% The assumption register. Four columns force the author to say where each
% number came from and which section tests it. If a row has no source, the
% number is a guess and should be labelled as one.
\begin{table}[H]
\caption{Assumption register. ``Tested in'' points at the sensitivity that
would expose the assumption.}
\small
\begin{tabularx}{\linewidth}{@{}L{4.4cm}L{3.0cm}YL{1.9cm}@{}}
\toprule
\sffamily\bfseries Assumption & \sffamily\bfseries Value used & \sffamily\bfseries Source & \sffamily\bfseries Tested in \\
\midrule
\ph{Discount rate} & \ph{5\% real} & \ph{course brief, \S2.3} & \ph{\S5.1} \\
\ph{Hydrogen reference price} & \ph{£7.80/kg} & \ph{policy strike, deflated to 2020 money} & \ph{\S5.1} \\
\ph{Battery round-trip efficiency} & \ph{90\%} & \ph{brief; 20--80\% usable window} & \ph{\S5.2} \\
\bottomrule
\end{tabularx}
\srcnote{Money is constant \ph{2020} GBP unless stated. Energy figures are
daily averages of metered half-hourly data.}
\end{table}

% ============================================================ section 3 ====
\section{Model}

\subsection{Decision variables}

\ph{One sentence per variable group, or a compact list. A reader must be able
to reconstruct the model from \S3.1--\S3.3 without the code.}

\subsection{Objective}

\begin{equation}
\max\ \Pi=\frac{1}{Y}\sum_{t}\left[\frac{p_t\,(q_t+d_t)}{1000}+\frac{v\,h_t}{\eta}\right]
-a_B\,n_B-a_H\,n_H
\label{eq:objective}
\end{equation}

\ph{Define every symbol immediately below the equation, in running text, not
in a symbol table the examiner has to hunt for. Say what the units are and
which term carries the money. Refer back to it as \eqref{eq:objective} --- a
number that is never cited in the text is a number nobody reads. A useful
test: give the magnitude of each term for the base case, so a reader can see
which one dominates and which one is a rounding error.}

% A report that quotes an equivalent annual value must show how the annual
% charge is built, otherwise the reader cannot reproduce a single row of the
% cash bridge.
\ph{If the report quotes equivalent annual values, show the annualisation
rather than describing it:}
\begin{equation}
a_j=F_j+\frac{i}{1-(1+i)^{-N_j}}\Bigl(K_j+\sum_{k\in R_j}\frac{C_{jk}}{(1+i)^{k}}\Bigr),
\qquad j\in\{B,H\}
\label{eq:annual}
\end{equation}
\ph{with the life, discount rate and renewal years named in the text, and at
least one worked value --- fixed cost plus capital plus renewals for the
asset that carries the decision.}

\subsection{Constraints}

\begin{align}
q_t+c_t+h_t+u_t &= G_t && \text{(daily energy balance)} \\
q_t+d_t &\le D_t && \text{(demand cap)} \\
0.096\,z_t \le h_t &\le 0.48\,z_t \le 0.48\,n_H && \text{(module limits)} \\
h_t &\le \eta\,L/10^{6} && \text{(offtake cap)} \\
\textstyle\sum_t d_t &= 0.9\sum_t c_t && \text{(annual battery balance)}
\end{align}

\ph{After the block, one sentence per constraint that needs justification,
and one sentence naming the relaxation. If a constraint is an upper bound on
value rather than a representation of physics, say so here, not in the
limitations.}

\subsection{Dispatch policies}

\ph{If the report compares an optimised allocation with a rule, define both
as maths, not as prose. Policy A: economic allocation. Policy B: current
generation serves current demand first, i.e.\ $q_t=\min(G_t,D_t)$. State which
one the client is permitted to run --- that is a commercial fact, not a
modelling choice.}

\subsection{Instance size and solver evidence}

% Solver evidence belongs in the body. Every cell must be a measurement or a
% literal N/A. A table with no failed cells is a table that was not tested.
\begin{table}[H]
\caption{Solver evidence for one design year. Times are wall-clock on
\ph{hardware}, default solver parameters except where noted.}
\small
\begin{tabularx}{\linewidth}{@{}L{4.4cm}R{1.4cm}R{1.4cm}R{2.1cm}R{1.3cm}Y@{}}
\toprule
\sffamily\bfseries Case & {\sffamily\bfseries Rows} & {\sffamily\bfseries Cols}
  & {\sffamily\bfseries Objective} & {\sffamily\bfseries Time} & \sffamily\bfseries Status \\
{\unitcell} & {\unitcell} & {\unitcell} & {\unitcell £m/yr} & {\unitcell s} & \\
\midrule
\ph{Baseline, policy A} & \num{2556} & \num{2192} & \num{17139.3} & 0.05 & optimal \\
\ph{20 MW electrolyser} & \num{2556} & \num{2192} & \num{17147.9} & 0.06 & optimal \\
\ph{50 t/day contract} & N/A & N/A & N/A & N/A & not run (licence limit) \\
\bottomrule
\end{tabularx}
\srcnote{\ph{Solver and version, licence class, gap settings, hardware, and the
one setting that changed the answer. Write N/A, never leave a cell blank.}}
\end{table}

% ============================================================ section 4 ====
\section{Results}

\subsection{Validation ladder}

% Same KPIs down four rows: baseline, naive rule, optimised, perfect
% foresight. The ladder is what proves the optimiser earns its keep.
\begin{table}[H]
\caption{Validation ladder: identical KPIs across four solution methods.}
\small
\begin{tabularx}{\linewidth}{@{}L{4.6cm}R{1.9cm}R{1.7cm}R{1.6cm}Y@{}}
\toprule
\sffamily\bfseries Method & {\sffamily\bfseries Value}
  & {\sffamily\bfseries Unused} & {\sffamily\bfseries Coverage} & \sffamily\bfseries What it shows \\
{\unitcell} & {\unitcell £m/yr} & {\unitcell GWh} & {\unitcell ratio} & \\
\midrule
\ph{No investment} & 0 & \num{549.9} & 0.875 & \ph{reference case} \\
\ph{Naive rule, same asset} & \num{-1.065} & \num{549.9} & 0.875 & \ph{a fixed rule destroys value} \\
\ph{Optimised, same asset} & 8.572 & \num{536.0} & 0.874 & \ph{value of dispatch} \\
\ph{Imperfect foresight} & N/A & N/A & N/A & \ph{outside Tier 1} \\
\bottomrule
\end{tabularx}
\srcnote{Increments are measured against the matching no-investment case, not
against each other.}
\end{table}

\subsection{Configuration comparison}

\begin{table}[H]
\caption{Annual increment by configuration, £m/yr, for the reference year.}
\small
\begin{tabular}{@{}L{5.6cm}R{2.2cm}R{2.2cm}R{2.2cm}R{2.2cm}@{}}
\toprule
\sffamily\bfseries Policy & {\sffamily\bfseries None} & {\sffamily\bfseries Battery}
 & {\sffamily\bfseries Electrolyser} & {\sffamily\bfseries Both} \\
\midrule
\ph{Policy A: economic} & 0 & 0 & \bfseries 8.57 & 8.57 \\
\ph{Policy B: supply first} & 0 & 0 & 0 & 0 \\
\bottomrule
\end{tabular}
\srcnote{Zero means the asset is not selected, not that the asset earns
nothing. Bold marks the recommended configuration.}
\end{table}

\subsection{Where the energy comes from}

% This subsection is the honesty test. Split the input into avoided waste and
% diverted sale, and give the attribution interval, not a single number.
\begin{figure}[H]
\centering
\figslot{3.9cm}{Figure slot 2 --- energy attribution, 11\,cm $\times$ 3.9\,cm\\
\footnotesize two-panel: input split by source; renewable share on surplus days}
\caption{\ph{Attribution of electrolyser input to avoided waste and diverted
sale, with the daily lower and upper bound on the renewable share.}}
\label{fig:attribution}
\figsrc{Source: \ph{results file, case id}.}
\end{figure}

\subsection{Selection}

\ph{Three short paragraphs, one per rejected alternative. For each: the
quantity that decides it, the threshold at which it would be selected, and
whether the threshold is inside the plausible range. If an alternative fails
for a reason the model cannot see, say which model in Tier 2 will see it.}

% ============================================================ section 5 ====
\section{Sensitivity, break-even and switching points}

% This section is not optional and is not an appendix. It is where a report
% stops describing a scenario and starts supporting a decision.

\subsection{One-at-a-time sensitivity}

\begin{table}[H]
\caption{Effect of single-parameter changes on the annual increment (£m/yr), capacity
held at the base case unless stated.}
\small
\begin{tabularx}{\linewidth}{@{}L{4.0cm}L{3.4cm}R{2.0cm}Y@{}}
\toprule
\sffamily\bfseries Parameter & \sffamily\bfseries Setting tested
  & {\sffamily\bfseries Increment} & \sffamily\bfseries Capacity changes? \\
\midrule
\ph{Electrolyser CAPEX} & \ph{£0.82m/MW, legacy} & 9.41 & no \\
\ph{Electrolyser CAPEX} & \ph{£1.65m/MW, high} & 8.33 & no \\
\ph{Discount rate} & \ph{0\% real} & 9.44 & no \\
\ph{Discount rate} & \ph{12\% real} & 6.99 & no \\
\ph{Hydrogen price} & \ph{£6.00/kg} & 3.33 & \ph{20 MW} \\
\ph{Hydrogen price} & \ph{£9.50/kg} & 13.88 & \ph{60 MW} \\
\bottomrule
\end{tabularx}
\srcnote{State the number of cases and how they were generated; a sensitivity
table with six rows and no description of the sweep is still a guess.}
\end{table}

\subsection{Break-even and switching points}

\begin{table}[H]
\caption{Thresholds at which the recommendation changes.}
\small
\begin{tabularx}{\linewidth}{@{}L{7.4cm}R{2.0cm}Y@{}}
\toprule
\sffamily\bfseries Threshold & {\sffamily\bfseries Value} & \sffamily\bfseries Comment \\
\midrule
\ph{Netback needed at 10 t/day} & 4.623 & \ph{£/kg; below this, do not build} \\
\ph{Offtake needed at £7.80/kg} & \num{2230} & \ph{kg/day} \\
\ph{Battery break-even capital cost} & \num{6990} & \ph{£/MWh on the low-FOM table} \\
\bottomrule
\end{tabularx}
\end{table}

\subsection{Critical map}

\ph{Two decision parameters on the axes, the decision boundary drawn, and the
base case marked. One figure replaces a page of two-way sensitivity prose.}

\subsection{What would falsify this}

\ph{Three to five bullets, each a measurable event: a contract term, a price,
a cost, or an obligation. ``If the buyer requires X, the investment is zero''
is falsifiable; ``results may vary'' is not.}

% ============================================================ section 6 ====
\section{Recommendation and implementation}

\begin{enumerate}
  \item \ph{Do this now --- with the number that justifies it.}
  \item \ph{Negotiate this first --- the term that gates the investment.}
  \item \ph{Stage the expansion --- what to build at which contract size.}
  \item \ph{Keep this alternative visible --- the obligation under which the
        answer reverses.}
\end{enumerate}

% ============================================================ section 7 ====
\section{Validation and limitations}

\subsection{Checks performed}

\ph{What was independently checked, how many checks, and what a check means
here (energy balance closure, cap compliance, recomputation of revenue from
the printed plan). Then the reproduction statement: which reference values
were reproduced, and the sentence that distinguishes arithmetic agreement
from endorsement of shared assumptions.}

\subsection{Quantified concessions}

% Every simplification gets a measured cost. An unquantified limitation is an
% admission; a quantified one is a design decision.
\begin{table}[H]
\caption{Approximations and their measured cost.}
\small
\begin{tabularx}{\linewidth}{@{}L{4.6cm}YL{3.4cm}@{}}
\toprule
\sffamily\bfseries Simplification & \sffamily\bfseries Cost, measured & \sffamily\bfseries Where it bites \\
\midrule
\ph{Daily time step} & \ph{no intraday arbitrage; battery value overstated} & \ph{\S4.5} \\
\ph{Annual battery balance} & \ph{relaxation, so an upper bound on value} & \ph{battery cases only} \\
\ph{Perfect foresight} & \ph{not quantified in Tier 1} & \ph{all cases} \\
\bottomrule
\end{tabularx}
\end{table}

\subsection{Possible objections}

\begin{caveatbox}[Weakest link, volunteered]
\ph{Name the parameter you trust least, say why it is still the right choice,
and give the number that would change your mind.}
\end{caveatbox}

% ============================================================ appendix =====
\section*{Reproducibility appendix}

\ph{Everything a second team needs to regenerate every number: data release
tag, package version, solver version and licence class, entry-point commands
in order, runtime, and the digest of the results file. A one-page appendix of
plain commands is worth more than a page of prose about methodology.}

{\small
\begin{tabularx}{\linewidth}{@{}L{6.2cm}Y@{}}
\toprule
\sffamily\bfseries Artefact & \sffamily\bfseries Command or digest \\
\midrule
\ph{Prepare inputs} & \texttt{\ph{python src/prepare\_data.py}} \\
\ph{Run all cases} & \texttt{\ph{python src/run\_cases.py}} \\
\ph{Verify} & \texttt{\ph{python src/verify.py}} \\
\ph{Results digest} & \texttt{\ph{sha256: 0000...}} \\
\bottomrule
\end{tabularx}}

% ------------------------------------------------------------ references ---
\section*{References}
{\small
\begin{enumerate}[label={[\arabic*]},leftmargin=2.1em]
  \item Group 15 data-repair release \ph{20261104-v2}, \ph{2,452 days}, accessed
        \ph{access date}; the repair log travels with the release.
  \item MMCS 2026/27 project brief, \ph{pp. 2--6}, and the teacher clarifications: asset
        modules, 10\,t/day offtake cap, 5\% real discount rate, £9.5/kg benchmark (2023 money).
  \item Office for National Statistics, Consumer Prices Index including owner occupiers'
        housing costs, series D7BT, monthly: the deflator to 2020 GBP.
  \item Cost sources with vintage: BEIS, \emph{Hydrogen Production Costs 2021} (2020 GBP);
        IEA, \emph{Global Hydrogen Review 2025}; Danish Energy Agency, \emph{Technology Data
        for Energy Storage}; \mbox{ITM POSEIDON} data sheet.
  \item DESNZ, \emph{HAR1 successful projects}, December 2023: weighted strike price
        £241/MWh (£175/MWh in 2012 prices). \ph{Name the independent implementation used for
        the cross-check: repository path, commit, date.}
\end{enumerate}}

% ============================================================== guide =====
\vspace{8pt}
\thispagestyle{plain}
\section*{Template guide --- delete before submission}

\subsection*{Page budget}

{\small
\begin{tabularx}{\linewidth}{@{}L{5.0cm}C{1.8cm}Y@{}}
\toprule
\sffamily\bfseries Section & \sffamily\bfseries Pages & \sffamily\bfseries What it must contain \\
\midrule
Title block, key figures, decision & 1 & four numbers and a reversible claim \\
Executive summary & 1 & numbers first, mechanism, the carrying assumption \\
Boundary and data & 1 & provenance table and assumption register \\
Model & 2 & objective, constraints, solver evidence with N/A cells \\
Results & 3 & validation ladder, configuration matrix, attribution figure \\
Sensitivity & 2 & one-at-a-time, break-even, critical map, falsifiers \\
Recommendation & 1 & numbered actions with numbers attached \\
Validation and limitations & 1 & checks, quantified concessions, weakest link \\
Appendix and references & 1--2 & commands, digests, sourced references \\
\bottomrule
\end{tabularx}}

\subsection*{Conventions this style enforces}

\begin{itemize}
  \item Figures and tables are numbered and referenced in text; captions sit
        above tables and below figures, and carry a source note.
  \item Units live in the column header, not in every cell. Numbers use a
        thousands separator and a fixed number of decimals per column.
  \item No vertical rules. Three horizontal rules per table, no more.
  \item A cell that was not measured reads \texttt{N/A}; a blank cell is a
        defect. In a numeric column write \texttt{\{N/A\}} so the alignment
        survives.
  \item Sensitivity, break-even and switching points have a fixed section of
        their own because they are the most common omission in past reports.
  \item Figure-to-table ratio of at least 1:1 in the results section; a
        report with sixteen tables and two figures reads as a spreadsheet.
\end{itemize}

\subsection*{Common defects this skeleton is designed to prevent}

\begin{enumerate}
  \item Opening with background instead of numbers.
  \item A recommendation with no threshold at which it reverses.
  \item Assumptions in prose, untraceable to a source.
  \item Solver claims with no instance size and no failed cells.
  \item A limitation section that is a list of apologies rather than measured
        costs.
  \item Money and physical quantities added inside one expression.
\end{enumerate}

\end{document}
